% ============================================
% LUC-17 — O Medo (edição revista R1075)
% ============================================
\fragtitle{luc}{L}{LUC}{17}{O Medo}
\markboth{LUC-17 --- O Medo}{}

\noindent \textit{Belo Horizonte --- 12 de abril de 2026, 3h14}

\noindent Lúcia acordou antes de olhar o relógio. Sabia que eram 3h14. Sabia pelo peso do ar, pelo modo como o silêncio do apartamento ficava mais espesso, como se alguém tivesse entrado e estivesse parado no corredor, esperando que ela se mexesse.

\noindent Ficou imóvel. O teto era uma superfície branca, lisa, sem nenhuma mancha. Contou até sessenta e dois. Depois levantou.

\noindent O diário estava na sala, dentro da caixa de metal, sobre a mesa. A chave estava num copo com água, ao lado. Lúcia tinha feito aquilo na noite anterior, depois de ler a última página. Era um gesto absurdo, ela sabia. Mas se acordasse com a chave seca na gaveta, saberia que precisava parar. Se acordasse com a água morna, saberia que algo tinha acontecido.

\noindent A água estava morna.

\noindent Ela não tocou na chave. Foi até a cozinha, bebeu um copo de água, olhou pela janela. A rua estava vazia. Um cachorro atravessou a calçada, sem pressa, e sumiu atrás de um carro. O olho direito dele brilhou quando passou sob o poste. Lúcia piscou. Quando abriu os olhos, o cachorro já não estava.

\noindent Voltou para a sala. Sentou no chão, de costas para o sofá, com a caixa de metal à frente. Não abriu. Ficou olhando para o cadeado, para a chave dentro do copo, para o reflexo da lâmpada na água.

\noindent Havia uma pasta sobre a mesa, com uma folha em branco e uma caneta. Lúcia pegou a caneta. Escreveu:

\noindent ``Mãe, se algo acontecer comigo, não leia o que está na caixa. Não importa o que encontrarem. Não importa o que eu tiver escrito. Queime.''

\noindent Olhou para a frase. A letra era a dela. Mas o ``M'' de ``Mãe'' estava inclinado para a esquerda, com um traço longo na perna direita. Ela não escrevia assim. Conhecia a própria letra desde a infância: redonda, apertada, sempre para a direita. Aquela letra era de outra mão.

\noindent Rasgou a folha. Rasgou em quatro, em oito, em dezesseis. Jogou os pedaços no lixo. Depois tirou os pedaços do lixo, colocou na pia, acendeu um fósforo. O papel queimou rápido, com uma chama azul no centro. As cinzas ficaram na pia, finas, escuras, com um cheiro doce que não era de papel queimado.

\noindent Lúcia abriu a torneira. Lavou as cinzas. Elas desceram pelo ralo, mas o cheiro ficou.

\noindent Ela foi até a porta. Verificou a fechadura. Estava trancada. Verificou a janela da sala. Trancada. A do quarto. Trancada. A da cozinha. Trancada. Voltou para a sala. Sentou de novo no chão. O apartamento era pequeno; ela conhecia cada ruído. O frigorífico. O cano da parede. O elevador no poço. Naquela noite, nenhum desses ruídos existia. O silêncio era total, como se o mundo tivesse sido coberto por uma camada de algodão.

\noindent Então ouviu a respiração.

\noindent A respiração era mais lenta, mais funda, vindo do corredor. Lúcia não se virou. Ficou olhando para a caixa de metal. A respiração parou. O silêncio voltou. Ela esperou. Nada.

\noindent ``Eu sei que você está aí'', ela disse.

\noindent A voz saiu baixa, firme, com uma calma que não sentia. O apartamento não respondeu.

\noindent ``Eu não vou fingir que não estou com medo'', ela continuou. ``Estou. Mas o medo não decide por mim. Quem decide sou eu.''

\noindent A água no copo tremeu. Não havia vento. Não havia vibração. A água tremeu como se alguém tivesse batido na mesa, muito de leve, do outro lado.

\noindent Lúcia se levantou. Pegou o copo. Bebeu a água. Não sentiu gosto de nada. Colocou o copo vazio na mesa. Pegou a chave. A chave estava quente, como se tivesse sido segurada por muito tempo. Ela a encaixou no cadeado. Abriu a caixa.

\noindent O diário estava lá, envolto no pano preto. Ela o tirou. Pesava mais do que devia. Abriu na última página. A letra era de Joaquim. A data era de 1979. No alto da página, escrito com outra tinta, havia uma frase que ela não tinha visto antes:

\noindent ``Você demorou.''

\noindent Lúcia fechou o diário. Ficou com ele no colo, olhando para a parede. O medo continuava ali, mas tinha mudado de lugar. Já não estava no corredor. Estava dentro dela, como um hóspede que não pede licença.

\noindent Ela pensou: \textit{eu ainda posso parar}. Podia devolver o diário, pedir afastamento, nunca mais abrir. Podia ligar para Ronaldo e dizer que desistia. Podia pegar o carro e dirigir até a casa da mãe. Podia fazer qualquer coisa, menos aquilo que sabia que ia fazer.

\noindent Abriu o diário de novo. Leu a primeira linha da última página. Depois a segunda. Depois a terceira. Quando chegou ao fim, já não havia mais medo. Havia uma espécie de fome.

\noindent\textit{\textbf{↪ Rotas:} \rota{LUC-13} \rotasep \rota{LUC-18} \rotasep \rota{CONT-01}}
